\section{Methodology}
\label{sec:method}

This section presents the Deep Conformal Alignment (DCA) framework for learning topology-preserving, angle-preserving transformations that normalize morphological variability in histopathology images. We first provide an overview of the key innovations (Section~\ref{subsec:overview}), then formalize the problem (Section~\ref{subsec:formulation}), detail the diffeomorphic transformation network (Section~\ref{subsec:diffeomorphic}), introduce the conformal regularization (Section~\ref{subsec:conformal}), and present the end-to-end learning objective (Section~\ref{subsec:objective}).

\subsection{Overview and Key Innovations}
\label{subsec:overview}

Figure~\ref{fig:architecture} illustrates the DCA architecture. Unlike standard CNNs that directly process raw images, DCA introduces a \textit{learnable geometric normalization module} that warps inputs into canonical representations before classification. Our framework makes three key technical contributions that distinguish it from prior work:

\textbf{(1) Diffeomorphic Spatial Transformer Network (D-STN):} While standard STNs~\cite{jaderberg2015spatial} learn unconstrained transformations that may introduce topological artifacts (e.g., folding, tearing), we parameterize deformations via velocity fields integrated through scaling-and-squaring~\cite{arsigny2006logdiff}, guaranteeing diffeomorphism---a critical requirement for biomedical imaging where tissue connectivity must be preserved.

\textbf{(2) Conformal Energy Regularization:} We introduce a differentiable loss derived from the Cauchy-Riemann equations in complex analysis, which constrains learned transformations to be approximately conformal (angle-preserving). This maintains local tissue structure integrity while normalizing global morphology. Our formulation connects to the Beltrami coefficient~\cite{beltrami_ref}, providing theoretical grounding in quasiconformal mapping theory.

\textbf{(3) End-to-End Joint Optimization:} Unlike prior geometric normalization methods~\cite{huang2024surface} that require expensive per-image optimization ($O(n^3)$ for $n$ pixels), DCA amortizes this cost through learning, enabling real-time inference ($O(n)$ forward pass) while adapting transformations to maximize task-specific performance.

\textbf{Comparison to Prior Work:} Table~\ref{tab:comparison_prior} contrasts DCA with related approaches. Standard STN learns unconstrained affine/thin-plate spline transformations without geometric constraints. Fixed conformal mapping~\cite{huang2024surface} enforces conformality but requires per-image optimization and cannot adapt to semantic content. DCA uniquely combines learnable transformations with explicit geometric constraints.

\begin{table}[t]
\centering
\caption{Comparison of Geometric Transformation Approaches}
\label{tab:comparison_prior}
\begin{tabular}{lccc}
\hline
\textbf{Method} & \textbf{Learnable} & \textbf{Conformal} & \textbf{Diffeomorphic} \\
\hline
Standard STN~\cite{jaderberg2015spatial} & \checkmark & $\times$ & $\times$ \\
Fixed CEM~\cite{huang2024surface} & $\times$ & \checkmark & \checkmark \\
\textbf{DCA (Ours)} & \checkmark & \checkmark & \checkmark \\
\hline
\end{tabular}
\end{table}

\subsection{Problem Formulation}
\label{subsec:formulation}

\subsubsection{Notation and Definitions}

Let $\Omega \subset \mathbb{R}^2$ denote the image domain (typically $[0,H] \times [0,W]$ for an $H \times W$ image) and $\mathcal{I} = L^2(\Omega, \mathbb{R}^3)$ the space of RGB images. Given a labeled dataset $\mathcal{D} = \{(I_i, y_i)\}_{i=1}^N$ where $I_i \in \mathcal{I}$ and $y_i \in \{1, \ldots, C\}$ denotes the tissue class, we seek to learn:

\begin{itemize}
    \item A \textit{transformation network} $\mathcal{T}_\psi: \mathcal{I} \rightarrow \mathfrak{X}(\Omega)$ that maps each image to a velocity field $v_i = \mathcal{T}_\psi(I_i)$, where $\mathfrak{X}(\Omega)$ denotes the Lie algebra of smooth vector fields on $\Omega$.
    \item A \textit{diffeomorphism} $\phi_i = \exp(v_i) \in \text{Diff}^+(\Omega)$, where $\text{Diff}^+(\Omega)$ denotes the group of orientation-preserving diffeomorphisms.
    \item A \textit{classifier} $f_\theta: \mathcal{I} \rightarrow \Delta^{C-1}$ that predicts class probabilities, where $\Delta^{C-1}$ is the $(C-1)$-dimensional probability simplex.
\end{itemize}

Table~\ref{tab:notation} summarizes key notation used throughout this section.

\begin{table}[t]
\centering
\caption{Mathematical Notation}
\label{tab:notation}
\begin{tabular}{ll}
\hline
\textbf{Symbol} & \textbf{Description} \\
\hline
$\Omega \subset \mathbb{R}^2$ & Image domain \\
$I: \Omega \rightarrow \mathbb{R}^3$ & RGB image \\
$v: \Omega \rightarrow \mathbb{R}^2$ & Velocity field \\
$\phi: \Omega \rightarrow \Omega$ & Diffeomorphism (deformation) \\
$J_\phi \in \mathbb{R}^{2 \times 2}$ & Jacobian matrix of $\phi$ \\
$\text{Diff}^+(\Omega)$ & Group of orientation-preserving diffeomorphisms \\
$\mathfrak{X}(\Omega)$ & Lie algebra of vector fields \\
$\mathcal{E}_{conf}(\phi)$ & Conformal energy \\
\hline
\end{tabular}
\end{table}

\subsubsection{Optimization Objective}

The transformed image is $\tilde{I}_i = I_i \circ \phi_i$, obtained by warping $I_i$ through $\phi_i$ via differentiable bilinear interpolation. We formulate the learning problem as:
\begin{equation}
\min_{\psi, \theta} \quad \mathbb{E}_{(I,y) \sim \mathcal{D}} \left[ \mathcal{L}_{cls}(f_\theta(\tilde{I}), y) + \lambda_{conf} \mathcal{E}_{conf}(\phi) + \lambda_{smooth} \mathcal{R}_{smooth}(\phi) \right]
\label{eq:objective}
\end{equation}
where:
\begin{itemize}
    \item $\mathcal{L}_{cls}$: cross-entropy classification loss
    \item $\mathcal{E}_{conf}(\phi)$: conformal energy measuring deviation from angle preservation
    \item $\mathcal{R}_{smooth}(\phi)$: smoothness regularizer preventing unrealistic deformations
    \item $\lambda_{conf}, \lambda_{smooth} > 0$: regularization weights
\end{itemize}

\textbf{Assumptions:} We make the following assumptions: (A1) The velocity field $v$ is sufficiently smooth ($v \in C^2(\Omega)$) to ensure $\phi = \exp(v)$ is a diffeomorphism; (A2) The optimal canonical representation exists and is approximately conformal; (A3) The classifier $f_\theta$ has sufficient capacity to learn discriminative features from normalized images.

\subsection{Diffeomorphic Spatial Transformer Network}
\label{subsec:diffeomorphic}

\subsubsection{Velocity Field Parameterization}

Directly parameterizing $\text{Diff}^+(\Omega)$ is challenging as it is an infinite-dimensional Lie group. We adopt the Lie algebra approach~\cite{arsigny2006logdiff,voxelmorph}: instead of predicting $\phi$ directly, we predict a stationary velocity field $v: \Omega \rightarrow \mathbb{R}^2$ and obtain $\phi$ via the exponential map.

\begin{definition}[Exponential Map]
\label{def:exp_map}
The exponential map $\exp: \mathfrak{X}(\Omega) \rightarrow \text{Diff}^+(\Omega)$ is defined as:
\begin{equation}
\phi = \exp(v) = \lim_{T \rightarrow \infty} \left( \text{Id} + \frac{v}{T} \right)^T
\end{equation}
where $\text{Id}$ denotes the identity transformation.
\end{definition}

Let $g_\psi: \mathcal{I} \rightarrow L^2(\Omega, \mathbb{R}^2)$ be a convolutional neural network (the \textit{localization network}). Given input $I$, we compute:
\begin{equation}
v = g_\psi(I)
\end{equation}

\textbf{Architecture Details:} We implement $g_\psi$ as a U-Net~\cite{ronneberger2015unet} with the following specifications:
\begin{itemize}
    \item \textbf{Encoder:} 4 downsampling blocks, each with two $3 \times 3$ convolutions + ReLU + $2 \times 2$ max pooling. Channel dimensions: $[3, 64, 128, 256, 512]$.
    \item \textbf{Decoder:} 4 upsampling blocks with skip connections, each with $2 \times 2$ transposed convolution + two $3 \times 3$ convolutions + ReLU.
    \item \textbf{Output Layer:} $1 \times 1$ convolution producing 2 channels (velocity field components $v_x, v_y$).
    \item \textbf{Initialization:} Final layer weights and biases initialized to zero, ensuring $v \approx 0$ at initialization (identity transformation).
\end{itemize}

\subsubsection{Scaling-and-Squaring Integration}

To compute $\phi = \exp(v)$, we use the scaling-and-squaring algorithm~\cite{arsigny2006logdiff}, which approximates the exponential map via repeated self-composition.

\begin{algorithm}[t]
\caption{Scaling-and-Squaring for Diffeomorphic Integration}
\label{alg:scaling_squaring}
\begin{algorithmic}[1]
\REQUIRE Velocity field $v \in \mathbb{R}^{H \times W \times 2}$, number of steps $T \in \mathbb{N}$
\ENSURE Deformation field $\phi \approx \exp(v)$
\STATE $\phi_0 \leftarrow v / 2^T$ \COMMENT{Scale velocity to small magnitude}
\FOR{$t = 1$ to $T$}
    \STATE $\phi_t(x) \leftarrow \phi_{t-1}(x) + \phi_{t-1}(x + \phi_{t-1}(x))$ \COMMENT{Self-composition}
    \STATE \hspace{1cm} \textit{// Implemented via bilinear interpolation}
\ENDFOR
\RETURN $\phi_T$
\end{algorithmic}
\end{algorithm}

The composition $\phi \circ \phi$ in Line 3 is computed via differentiable bilinear interpolation: given $\phi(x) = (u(x), v(x))$, we evaluate $\phi(\phi(x)) = \phi(x + \phi(x))$ by sampling $\phi$ at location $x + \phi(x)$.

\begin{proposition}[Diffeomorphism Guarantee]
\label{prop:diffeomorphism}
If $\|v\|_{L^\infty} < \epsilon$ for sufficiently small $\epsilon > 0$, then $\phi = \exp(v)$ obtained via Algorithm~\ref{alg:scaling_squaring} with $T \geq 7$ is a diffeomorphism with probability $> 1 - \delta$ for arbitrarily small $\delta > 0$.
\end{proposition}

\begin{proof}[Proof Sketch]
The result follows from~\cite{arsigny2006logdiff}: for small $v$, the exponential map $\exp: \mathfrak{X}(\Omega) \rightarrow \text{Diff}^+(\Omega)$ is a local diffeomorphism. The scaling-and-squaring approximation converges exponentially fast: $\|\phi_T - \exp(v)\| = O(2^{-T})$. For $T=7$, the approximation error is $< 10^{-3}$ in practice. The condition $\|v\|_{L^\infty} < \epsilon$ ensures the Jacobian determinant $\det(J_\phi) > 0$ everywhere, guaranteeing orientation preservation.
\end{proof}

\textbf{Computational Complexity:} Algorithm~\ref{alg:scaling_squaring} requires $T$ self-compositions, each involving bilinear interpolation over $H \times W$ pixels. The total complexity is $O(THW)$, which is linear in image size.

\subsection{Conformal Regularization via Differential Geometry}
\label{subsec:conformal}

\subsubsection{Conformal Mappings and Cauchy-Riemann Equations}

A smooth map $\phi: \Omega \rightarrow \Omega$ is \textit{conformal} if it preserves angles locally. In 2D, conformality is characterized by the Cauchy-Riemann equations from complex analysis.

\begin{definition}[Conformal Map]
\label{def:conformal}
Let $\phi(x,y) = (u(x,y), v(x,y))$ be a smooth map. The Jacobian matrix is:
\begin{equation}
J_\phi = \begin{bmatrix} \frac{\partial u}{\partial x} & \frac{\partial u}{\partial y} \\ \frac{\partial v}{\partial x} & \frac{\partial v}{\partial y} \end{bmatrix}
\end{equation}
The map $\phi$ is conformal if and only if $J_\phi$ is a scaled rotation:
\begin{equation}
J_\phi = \lambda(x,y) R(\theta(x,y)) = \lambda(x,y) \begin{bmatrix} \cos\theta & -\sin\theta \\ \sin\theta & \cos\theta \end{bmatrix}
\label{eq:conformal_jacobian}
\end{equation}
for some scaling factor $\lambda > 0$ and rotation angle $\theta$. This is equivalent to the Cauchy-Riemann equations:
\begin{equation}
\frac{\partial u}{\partial x} = \frac{\partial v}{\partial y}, \quad \frac{\partial u}{\partial y} = -\frac{\partial v}{\partial x}
\label{eq:cauchy_riemann}
\end{equation}
\end{definition}

\textbf{Geometric Interpretation:} Conformal maps preserve angles but not necessarily lengths. For histopathology, this means local tissue structures (e.g., cell arrangements, gland orientations) maintain their angular relationships while global morphology is normalized.

\subsubsection{Differentiable Conformal Energy}

Since our transformation is $F(x) = x + \phi(x)$ (displacement formulation), we measure conformality of $F$. The Jacobian is $J_F = I + J_\phi$. We define the conformal energy as the $L^2$ deviation from Eq.~\eqref{eq:cauchy_riemann}:

\begin{equation}
\mathcal{E}_{conf}(\phi) = \frac{1}{|\Omega|} \int_\Omega \left[ \left( \frac{\partial u}{\partial x} - \frac{\partial v}{\partial y} \right)^2 + \left( \frac{\partial u}{\partial y} + \frac{\partial v}{\partial x} \right)^2 \right] dx
\label{eq:conformal_energy}
\end{equation}

This can be equivalently written using the Frobenius norm:
\begin{equation}
\mathcal{E}_{conf}(\phi) = \frac{1}{|\Omega|} \int_\Omega \left\| J_\phi - \frac{\text{tr}(J_\phi)}{2} I \right\|_F^2 dx
\end{equation}
where $\text{tr}(J_\phi) = \frac{\partial u}{\partial x} + \frac{\partial v}{\partial y}$ is the trace.

\textbf{Discretization:} In practice, we discretize using central finite differences:
\begin{align}
\frac{\partial u}{\partial x}(i,j) &\approx \frac{u(i+1,j) - u(i-1,j)}{2\Delta x} \\
\frac{\partial u}{\partial y}(i,j) &\approx \frac{u(i,j+1) - u(i,j-1)}{2\Delta y}
\end{align}
and similarly for $v$. The discrete conformal energy is:
\begin{equation}
\mathcal{E}_{conf}(\phi) = \frac{1}{HW} \sum_{i,j} \left[ \left( \nabla_x u_{ij} - \nabla_y v_{ij} \right)^2 + \left( \nabla_y u_{ij} + \nabla_x v_{ij} \right)^2 \right]
\label{eq:discrete_conformal}
\end{equation}

\subsubsection{Connection to Beltrami Coefficient}

The Beltrami coefficient $\mu: \Omega \rightarrow \mathbb{C}$ is a complex-valued measure of quasiconformality defined as:
\begin{equation}
\mu = \frac{\partial_{\bar{z}} \phi}{\partial_z \phi}
\end{equation}
where $z = x + iy$ and $\partial_z = \frac{1}{2}(\partial_x - i\partial_y)$, $\partial_{\bar{z}} = \frac{1}{2}(\partial_x + i\partial_y)$.

\begin{proposition}[Conformality and Beltrami Coefficient]
\label{prop:beltrami}
A map $\phi$ is conformal if and only if $\mu = 0$ almost everywhere. Furthermore:
\begin{equation}
\mathcal{E}_{conf}(\phi) = \frac{4}{|\Omega|} \int_\Omega |\mu|^2 \left| \frac{\partial \phi}{\partial z} \right|^2 dx
\end{equation}
\end{proposition}

\begin{proof}
From the definition of $\mu$ and the Cauchy-Riemann equations, $\mu = 0 \Leftrightarrow \partial_{\bar{z}} \phi = 0 \Leftrightarrow$ Eq.~\eqref{eq:cauchy_riemann} holds. The energy equivalence follows from direct computation using $\partial_z \phi = \frac{1}{2}[(\partial_x u - \partial_y v) + i(\partial_x v + \partial_y u)]$.
\end{proof}

This connection provides theoretical grounding: minimizing $\mathcal{E}_{conf}$ is equivalent to minimizing the $L^2$ norm of the Beltrami coefficient, which measures deviation from conformality in quasiconformal mapping theory~\cite{beltrami_ref}.

\subsection{End-to-End Learning Objective}
\label{subsec:objective}

\subsubsection{Complete Loss Function}

The full training objective combines three terms:
\begin{equation}
\mathcal{L}(\psi, \theta; I, y) = \mathcal{L}_{cls}(f_\theta(\tilde{I}), y) + \lambda_{conf} \mathcal{E}_{conf}(\phi) + \lambda_{smooth} \mathcal{R}_{smooth}(\phi)
\label{eq:full_loss}
\end{equation}
where:
\begin{itemize}
    \item \textbf{Classification Loss:} $\mathcal{L}_{cls} = -\sum_{c=1}^C y_c \log f_\theta(\tilde{I})_c$ (cross-entropy)
    \item \textbf{Conformal Energy:} $\mathcal{E}_{conf}(\phi)$ from Eq.~\eqref{eq:discrete_conformal}
    \item \textbf{Smoothness Regularizer:} $\mathcal{R}_{smooth}(\phi) = \frac{1}{|\Omega|} \sum_{x \in \Omega} \|\nabla \phi(x)\|_F^2$ (total variation)
\end{itemize}

\textbf{Hyperparameter Selection:} We set $\lambda_{conf} = 1.0$ and $\lambda_{smooth} = 0.1$ based on grid search over $\{0.1, 0.5, 1.0, 5.0\} \times \{0.01, 0.1, 1.0\}$ on the validation set. The conformal weight $\lambda_{conf}$ controls the trade-off between task performance and geometric constraint satisfaction.

\subsubsection{Optimization Algorithm}

We optimize Eq.~\eqref{eq:full_loss} using Adam~\cite{kingma2014adam} with the following schedule:

\begin{algorithm}[t]
\caption{DCA Training Procedure}
\label{alg:training}
\begin{algorithmic}[1]
\REQUIRE Dataset $\mathcal{D} = \{(I_i, y_i)\}_{i=1}^N$, batch size $B$, epochs $E$, learning rate $\alpha_0$
\ENSURE Trained parameters $\psi^*, \theta^*$
\STATE Initialize $\psi, \theta$ randomly (except final layer of $g_\psi$ initialized to zero)
\STATE Set learning rate $\alpha \leftarrow \alpha_0 = 10^{-3}$
\FOR{epoch $e = 1$ to $E$}
    \FOR{each mini-batch $\{(I_b, y_b)\}_{b=1}^B \sim \mathcal{D}$}
        \STATE $v_b \leftarrow g_\psi(I_b)$ \COMMENT{Predict velocity fields}
        \STATE $\phi_b \leftarrow \text{ScalingSquaring}(v_b, T=7)$ \COMMENT{Algorithm~\ref{alg:scaling_squaring}}
        \STATE $\tilde{I}_b \leftarrow I_b \circ \phi_b$ \COMMENT{Warp images via bilinear interpolation}
        \STATE $\hat{y}_b \leftarrow f_\theta(\tilde{I}_b)$ \COMMENT{Classify transformed images}
        \STATE Compute $\mathcal{L}(\psi, \theta; I_b, y_b)$ via Eq.~\eqref{eq:full_loss}
        \STATE $(\psi, \theta) \leftarrow \text{Adam}(\nabla_{\psi,\theta} \mathcal{L}, \alpha)$ \COMMENT{Update parameters}
    \ENDFOR
    \IF{$e \mod 10 = 0$}
        \STATE $\alpha \leftarrow \alpha \times 0.95$ \COMMENT{Learning rate decay}
    \ENDIF
\ENDFOR
\RETURN $\psi^*, \theta^*$
\end{algorithmic}
\end{algorithm}

\subsubsection{Computational Complexity Analysis}

Let $H \times W$ denote image resolution, $D$ the depth of $g_\psi$, $D'$ the depth of $f_\theta$, and $T=7$ the scaling-and-squaring steps. The per-image computational cost is:
\begin{equation}
\text{Cost} = O(HWD) + O(THW) + O(HWD') = O(HW(D + D' + T))
\end{equation}

\textbf{Training vs. Inference:} During training, we compute gradients through all components, requiring $O(HW(D + D' + T))$ per image. During inference, we perform a single forward pass with the same complexity. Critically, unlike fixed conformal mapping methods~\cite{huang2024surface} that require $O(n^3)$ per-image optimization (where $n = HW$), DCA amortizes this cost during training, enabling real-time deployment.

\subsubsection{Implementation Details}

\begin{itemize}
    \item \textbf{Framework:} PyTorch 1.12 with CUDA 11.3
    \item \textbf{Hardware:} NVIDIA A100 GPU (40GB memory)
    \item \textbf{Batch Size:} 16 (limited by GPU memory for $224 \times 224$ images)
    \item \textbf{Training Time:} $\sim$2 hours for 10 epochs on 5,026 training images
    \item \textbf{Inference Time:} $\sim$15ms per image (including transformation + classification)
    \item \textbf{Classifier Architecture:} ResNet-18~\cite{he2016resnet} pretrained on ImageNet, fine-tuned end-to-end
\end{itemize}

\textbf{Gradient Flow:} All components (velocity field prediction, scaling-and-squaring, bilinear sampling, classification) are differentiable, enabling end-to-end backpropagation. The gradient of the conformal energy w.r.t. $\phi$ is computed via automatic differentiation of Eq.~\eqref{eq:discrete_conformal}.
